% Mathematical TeX extracted from the cited web pages, reformatted by the assistant. Not the native full chapter.
Before we state the result we need some notation. Let $R$ be a ring. Recall that $R[T_0,\ldots,T_n]$ is a graded $R$-algebra where each $T_i$ is homogeneous of degree $1$. Denote $(R[T_0,\ldots,T_n])_d$ the degree $d$ summand. It is a finite free $R$-module of rank $\binom{n+d}{d}$ when $d\geq0$ and zero else. It has a basis consisting of monomials $T_0^{e_0}\ldots T_n^{e_n}$ with $\sum e_i=d$.

We will also use the following notation: $R[\frac1{T_0},\ldots,\frac1{T_n}]$ denotes the $\Z$-graded ring with $\frac1{T_i}$ in degree $-1$. In particular the $\Z$-graded $R[\frac1{T_0},\ldots,\frac1{T_n}]$ module
\[
\frac1{T_0\ldots T_n}R[\tfrac1{T_0},\ldots,\tfrac1{T_n}]
\]
which shows up in the statement below is zero in degrees $\geq-n$, is free on the generator $\frac1{T_0\ldots T_n}$ in degree $-n-1$ and is free of rank $(-1)^n\binom{n+d}{d}$ for $d\leq-n-1$.

\begin{lemma}[30.8.1]
Let $R$ be a ring. Let $n\geq0$ be an integer. We have
\[
H^q(\mathbf P^n,\mathcal O_{\mathbf P_R^n}(d))=
\begin{cases}
(R[T_0,\ldots,T_n])_d& q=0,\ d\geq0,\\
\left(\frac1{T_0\ldots T_n}R[\frac1{T_0},\ldots,\frac1{T_n}]\right)_d& q=n,\ d<0,\\
0& q,n,d\text{ not as above}
\end{cases}
\]
as $R$-modules.
\end{lemma}
\begin{proof}
We will use the standard affine open covering
\[
\mathcal U:\mathbf P_R^n=\bigcup_{i=0}^nD_+(T_i)
\]
to compute the cohomology using the \v Cech complex. This is permissible by Lemma 30.2.6 since any intersection of finitely many affine $D_+(T_i)$ is also a standard affine open (see Constructions, Section 27.8). In fact, we can use the alternating or ordered \v Cech complex according to Cohomology, Lemmas 20.23.3 and 20.23.6.

The ordering we will use on $\{0,\ldots,n\}$ is the usual one. Hence the complex we are looking at has terms
\[
\check{\mathcal C}_{ord}^p(\mathcal U,\mathcal O_{\mathbf P_R}(d))
=\bigoplus_{i_0<\ldots<i_p}(R[T_0,\ldots,T_n,\tfrac1{T_{i_0}\ldots T_{i_p}}])_d.
\]
Moreover, the maps are given by the usual formula
\[
d(s)_{i_0\ldots i_{p+1}}=\sum_{j=0}^{p+1}(-1)^js_{i_0\ldots\widehat{i_j}\ldots i_{p+1}}
\]
see Cohomology, Section 20.23. Note that each term of this complex has a natural $\Z^{n+1}$-grading. Namely, we get this by declaring a monomial $T_0^{e_0}\ldots T_n^{e_n}$ to be homogeneous with weight $(e_0,\ldots,e_n)\in\Z^{n+1}$. It is clear that the differential given above respects the grading. In a formula we have
\[
\check{\mathcal C}_{ord}^{\bullet}(\mathcal U,\mathcal O_{\mathbf P_R}(d))
=\bigoplus_{\vec e\in\Z^{n+1}}\check{\mathcal C}^{\bullet}(\vec e)
\]
where not all summands on the right hand side occur (see below). Hence in order to compute the cohomology modules of the complex it suffices to compute the cohomology of the graded pieces and take the direct sum at the end.

Fix $\vec e=(e_0,\ldots,e_n)\in\Z^{n+1}$. In order for this weight to occur in the complex above we need to assume $e_0+\ldots+e_n=d$ (if not then it occurs for a different twist of the structure sheaf of course). Assuming this, set
\[
NEG(\vec e)=\{i\in\{0,\ldots,n\}\mid e_i<0\}.
\]
With this notation the weight $\vec e$ summand $\check{\mathcal C}^{\bullet}(\vec e)$ of the \v Cech complex above has the following terms
\[
\check{\mathcal C}^{p}(\vec e)
=\bigoplus_{\substack{i_0<\ldots<i_p\\NEG(\vec e)\subset\{i_0,\ldots,i_p\}}}
R\cdot T_0^{e_0}\ldots T_n^{e_n}.
\]
In other words, the terms corresponding to $i_0<\ldots<i_p$ such that $NEG(\vec e)$ is not contained in $\{i_0\ldots i_p\}$ are zero. The differential of the complex $\check{\mathcal C}^{\bullet}(\vec e)$ is still given by the exact same formula as above.

Suppose that $NEG(\vec e)=\{0,\ldots,n\}$, i.e., that all exponents $e_i$ are negative. In this case the complex $\check{\mathcal C}^{\bullet}(\vec e)$ has only one term, namely $\check{\mathcal C}^{n}(\vec e)=R\cdot\frac1{T_0^{-e_0}\ldots T_n^{-e_n}}$. Hence in this case
\[
H^q(\check{\mathcal C}^{\bullet}(\vec e))=
\begin{cases}R\cdot\frac1{T_0^{-e_0}\ldots T_n^{-e_n}}&q=n,\\0&\text{else}.\end{cases}
\]
The direct sum of all of these terms clearly gives the value
\[
\left(\frac1{T_0\ldots T_n}R[\tfrac1{T_0},\ldots,\tfrac1{T_n}]\right)_d
\]
in degree $n$ as stated in the lemma. Moreover these terms do not contribute to cohomology in other degrees (also in accordance with the statement of the lemma).

Assume $NEG(\vec e)=\emptyset$. In this case the complex $\check{\mathcal C}^{\bullet}(\vec e)$ has a summand $R$ corresponding to all $i_0<\ldots<i_p$. Let us compare the complex $\check{\mathcal C}^{\bullet}(\vec e)$ to another complex. Namely, consider the affine open covering
\[
\mathcal V:\Spec(R)=\bigcup_{i\in\{0,\ldots,n\}}V_i
\]
where $V_i=\Spec(R)$ for all $i$. Consider the alternating \v Cech complex
\[
\check{\mathcal C}_{ord}^{\bullet}(\mathcal V,\mathcal O_{\Spec(R)}).
\]
By the same reasoning as above this computes the cohomology of the structure sheaf on $\Spec(R)$. Hence we see that $H^p(\check{\mathcal C}_{ord}^{\bullet}(\mathcal V,\mathcal O_{\Spec(R)}))=R$ if $p=0$ and is $0$ whenever $p>0$. For these facts, see Lemma 30.2.1 and its proof.

Note that also $\check{\mathcal C}_{ord}^{\bullet}(\mathcal V,\mathcal O_{\Spec(R)})$ has a summand $R$ for every $i_0<\ldots<i_p$ and has exactly the same differential as $\check{\mathcal C}^{\bullet}(\vec e)$. In other words these complexes are isomorphic complexes and hence have the same cohomology. We conclude that
\[
H^q(\check{\mathcal C}^{\bullet}(\vec e))=
\begin{cases}R\cdot T_0^{e_0}\ldots T_n^{e_n}&q=0,\\0&\text{else}\end{cases}
\]
in the case that $NEG(\vec e)=\emptyset$. The direct sum of all of these terms clearly gives the value $(R[T_0,\ldots,T_n])_d$ in degree $0$ as stated in the lemma. Moreover these terms do not contribute to cohomology in other degrees (also in accordance with the statement of the lemma).

To finish the proof of the lemma we have to show that the complexes $\check{\mathcal C}^{\bullet}(\vec e)$ are acyclic when $NEG(\vec e)$ is neither empty nor equal to $\{0,\ldots,n\}$. Pick an index $i_{\mathrm{fix}}\notin NEG(\vec e)$ (such an index exists). Consider the map
\[
h:\check{\mathcal C}^{p+1}(\vec e)\to\check{\mathcal C}^{p}(\vec e)
\]
given by the rule that for $i_0<\ldots<i_p$ we have
\[
h(s)_{i_0\ldots i_p}=\begin{cases}
0&p\notin\{0,\ldots,n-1\},\\
0&i_{\mathrm{fix}}\in\{i_0,\ldots,i_p\},\\
s_{i_{\mathrm{fix}}i_0\ldots i_p}&i_{\mathrm{fix}}<i_0,\\
(-1)^as_{i_0\ldots i_{a-1}i_{\mathrm{fix}}i_a\ldots i_p}&i_{a-1}<i_{\mathrm{fix}}<i_a,\\
(-1)^{p+1}s_{i_0\ldots i_pi_{\mathrm{fix}}}&i_p<i_{\mathrm{fix}}.
\end{cases}
\]
Please compare with the proof of Lemma 30.2.1. This makes sense because we have
\[
NEG(\vec e)\subset\{i_0,\ldots,i_p\}\ \Longleftrightarrow\
NEG(\vec e)\subset\{i_{\mathrm{fix}},i_0,\ldots,i_p\}.
\]
The exact same (combinatorial) computation as in the proof of Lemma 30.2.1 shows that
\[
(hd+dh)(s)_{i_0\ldots i_p}=s_{i_0\ldots i_p}.
\]
Hence we see that the identity map of the complex $\check{\mathcal C}^{\bullet}(\vec e)$ is homotopic to zero which implies that it is acyclic.
\end{proof}

